%% kapa: exam=1 round=36 year=2025 period=1 form=A subject=civil_law subject_ko=민법 seq=15 number=15
%% source: Q-Net 감정평가사 기출문제 (cst00309) artlSeq=5231525 file=raw/1008/5231525/2025년 제36회 감정평가사 1차시험 1교시 문제지 원본.pdf
%% conversion: kordoc 4.24.51 -> tools/gen1.py (auto); review=proofread-sonnet; flags=none
\begin{question}{15}
甲은 토지거래허가구역 내에 있는 자신의 X토지에 관하여 乙과 매매계약을 체결하고, 일정 기간 내에 토지거래허가를 받기로 합의하였다. 이에 관한 설명으로 옳지 않은 것은? (다툼이 있으면 판례에 따름)
\begin{choices}
  \item 甲과 乙 쌍방이 토지거래허가신청을 하지 않기로 하는 의사를 명백히 표시한 경우, 매매계약은 확정적으로 무효가 된다.
  \item 甲이 허가신청절차에 협력하지 않는다면 乙은 甲에 대하여 협력의무의 이행을 소구할 수 있다.
  \item 토지거래허가가 있기 전에는 乙은 매매계약에 따른 대금지급의무가 없다.
  \item 甲과 乙이 토지거래허가를 받으면 매매계약은 특별한 사정이 없는 한 허가를 받은 때부터 유효가 된다.
  \item 만약 丙이 乙로부터 X토지를 매수한 후 자신과 甲을 당사자로 하는 토지거래허가를 받아 甲으로부터 소유권이전등기를 경료받았다면 그 등기는 무효이다.
\end{choices}
\end{question}
\answer{4}
